% ============================================================================
% REMERCIEMENTS (Acrostiche M-E-R-C-I)
% Fichier : preambule/remerciements.tex
% ============================================================================

\chapter*{REMERCIEMENTS}
\markboth{Remerciements}{}
\addcontentsline{toc}{chapter}{Remerciements}
\thispagestyle{preliminaire}
\begin{sloppypar}
\initial{M}es remerciements les plus sinc\`eres s'adressent premi\`erement \`a mes directeurs de th\`ese, le Professeur \textbf{Thomas DJOTIO NDI\'E} et le Docteur \textbf{Flavien Serge MANI ONANA}, pour leur encadrement rigoureux et leurs expertises combin\'ees. Je tiens \'egalement \`a exprimer ma profonde gratitude aux membres du jury, xxxxxxxxxxxxxxxxxxxxxxxxxxxxxxxxxxxxxxxxxxxxxxxxxxxxxxxxxxxxxxxxxxxxxxxxxxxx xxxxxxxxxxxxxxxxxxxxxxxxxxxxxxxxxxxxxxxxxxxxxxxx, qui ont accept\'e d'\'evaluer ce travail. Dans la même lanc\'ee, Professeurs \textbf{Georges KOUAMOU, Achille MELINGUI} et \textbf{Isidore NGONGO} de
l'\'Ecole Nationale Supérieure Polytechnique de Yaoundé I, merci pour vos remarques constructives lors du jury d'audition. Elles ont \'et\'e essentielles pour porter ce travail \`a ce niveau d'exigence.
\end{sloppypar}
\vspace{0.5cm}

\initial{E}nvers les \textbf{Services du Contr\^ole Sup\'erieur de l'\'Etat (CONSUPE)}, mon employeur actuel, j'exprime toute ma reconnaissance pour la bienveillance qui m'a permis de concilier responsabilit\'es administratives et exigences doctorales. Ma gratitude s'\'etend aux magistrats civils et militaires avec qui j'ai eu l'honneur d'\'echanger en ma qualit\'e d'Expert Judiciaire. Il en va de m\^eme pour mes \'echanges nourris avec certains responsables du Minist\`ere des Finances et de l'H\^opital G\'en\'eral de Garoua. Ces interactions ont ancr\'e ma r\'eflexion sur la cryptographie institutionnelle dans la r\'ealit\'e juridique et les enjeux de souverainet\'e num\'erique.
\vspace{0.5cm}
\initial{R}echercher l'excellence m'a conduit \`a m'impliquer activement dans la communaut\'e scientifique internationale, notamment au sein de l'\textbf{IACR}, de l'\textbf{IEEE} et de l'\textbf{ISACA} dont je suis membre. Mon appartenance \`a ces organisations refl\`ete ma volont\'e de lier la recherche acad\'emique aux standards professionnels de la s\'ecurit\'e de l'information. Je remercie \'egalement tous mes co-auteurs. Puisse notre collaboration continuer. \`A tous les contributeurs de la communaut\'e \textbf{Lean4} pour leur aide pr\'ecieuse dans la formalisation de mes preuves de s\'ecurit\'e, je dis simplement \textit{you rock!}
\vspace{0.5cm}
\initial{C}e travail est aussi le fruit de l'\'emulation v\'ecue au sein du \textbf{Laboratoire d'Ing\'enierie Math\'ematique et des Syst\`emes d'Information (LIMSI)}. Je pense aux encouragements du Professeur \textbf{Thomas BOUETOU}, aux discussions chaleureuses avec mes coll\`egues doctorants, aux critiques constructives des enseignants du D\'epartement de G\'enie Informatique pendant les diff\'erentes rencontres doctorales, mais également \`a mes \'etudiants de diverses promotions (\textbf{GI} et \textbf{HN}). Savoir que certains d'entre vous me consid\`erent comme un mod\`ele est un honneur qui m'engage quotidiennement.
\vspace{0.5cm}
\initial{I}ndispensable a \'et\'e le soutien de mon cercle le plus cher, ma famille. \`A mon \'epouse Elisabeth, pour sa patience~; \`a mes fils David et Axel~; \`a mes filles Sheila et Bissa\"i, \textbf{les MINKA} qui m'ont souvent vu inaccessible \`a cause de ce projet~; et \`a mes parents, \textbf{les Nguidjo\"i}, dont la fiert\'e est ma plus belle r\'ecompense. \textbf{Que ce travail ouvre grand les portes du doctorat pour notre lign\'ee}.
\vspace{0.4cm}
\begin{flushright}
\textit{\`A tous et \`a chacun, du plus profond de mon \^etre, \textbf{MERCI !}}
\end{flushright}